% FIXME: "O foolish Galatians" when DAN 1871 is complete.

\section{Tenth Sunday after Trinity: To whom should we go?}


\begin{quote}

John 6:66–71. From that time many of his disciples turned back and no longer walked with him. Then Jesus said to the twelve: Will you also go away? Simon Peter answered him: Lord, to whom should we go? You have the words of eternal life, and we have believed and known that you are the Christ, the Son of the living God. Jesus answered them: Have I not chosen you twelve? And one of you is a devil. But he spoke of Judas Iscariot, Simon’s son; for he was the one who would betray him, and he was one of the twelve.

\end{quote}

\bigskip


In the sixth chapter of John we encounter a great and decisive turning-point, both in the relation of the people and in that of the larger circle of disciples to Jesus. This chapter may be read again and again with blessing beyond measure by all—unbelievers and apostates, weak Christians and strong alike—and in it each may find both relief for wounds and remedy for hard hearts. In it, every one may both see his own deceitful heart portrayed and find the Gospel’s bread of life.

Jesus had fed five thousand. Then they all became so enthusiastic for him that they even "would take him by force to make him king." They could not rest till they saw and heard him again—and were fed by him. They were willing to go however far it might be in order to find him again.

But have you ever been in such a state of mind?

Their enthusiasm became, if possible, even greater when they found him again in Capernaum and learned that he must have gone there across the sea. They pressed in upon him and surely thought that they would love him and cling to him all their lives. Here, they thought, was the king for Israel!

But here their enthusiasm suffered its first shock. The Lord, as it were, cuts into their hearts and lays bare before all what really lay behind their admiration and their surging love for him: "You seek me because you ate of the loaves and were filled!" But that the heart’s deepest thoughts and intentions should thus be revealed, they could not bear it.

But can you endure it?

Jesus’ admonition not to labor for the food which perishes, but for that which endures unto eternal life, was not to their liking, and they began—as so many have done and will do—they began wrangling about religion—about Moses and manna—and to demand signs. When Jesus explained to them that Moses gave no bread from heaven, but that the Son of Man himself was the true manna, the Bread of Life, who had come down from heaven and on the last day would raise up those who believed in him instead of asking for some other sign, then they seized upon a welcome occasion to murmur against Jesus and themselves reveal the true state of their hearts. "Is he not Joseph’s son, whose father and mother we know? and he has come down from heaven?" they said.

Now it had come to a decision with them—and a sorrowful one. They forgot both the feeding and the walking on the sea and their enthusiasm. And when Jesus further unfolded to them that Moses and the manna were food for dying men, and that they had to eat Jesus’ flesh and drink his blood in order to share his own relation to the living Father and receive eternal life, then they thought they had full right to regard him as a Samaritan and a madman. Many, even of his disciples, murmured and said: "This is a hard saying; who can hear it?"

And thus it came to pass that many of his disciples went back and no longer walked with Jesus.

But has the same thing happened to you?

To how many does Paul’s heart-rending cry apply, when he says: "O foolish Galatians, who has bewitched you, that you should not obey the truth, you before whose eyes Jesus Christ was portrayed as crucified among you! Are you so foolish? Having begun in the Spirit, will you now be made complete in the flesh?"

But do you not see Jesus’ tears, as he weeps over Jerusalem and says: "Jerusalem! how often I would have gathered you as a hen gathers her chicks under her wings, and you would not. O if only you had known, even on this your day, what makes for your peace! But now it is hidden from your eyes."

Those tears he still sheds today over you who no longer walk with him.

Will you also leave him, you twelve?

"To whom should we go?"

I sighed heavily under the burden of sin and Satan’s bondage, and your word concerning sin, O Lord, was like dagger-thrusts in my conscience, so that my flesh rebelled, and I was near to hardening myself against your truth.

But you did not let me go, you Shepherd of Israel; you had engraved me upon your hands, and when you led me to the foot of the cross, and in faith I was given to look up to my Savior’s thorn-crowned head and knew that the burden of sin was taken away by his blood, and I was free and saved—O what a glad and blessed hour that was! You said: "Take heart; your sins are forgiven you!" and I, I would never part from my Savior, in life and in death.

"To whom then shall I go?"

You have the words of eternal life. By your good Holy Spirit you are always near me in your word and in your body; when I hunger, I go to you and am satisfied; when I thirst, I draw from Siloam’s softly flowing spring and am refreshed. You lead me into green pastures, you surround me with delight, you have taught me a new song, a blessed song of joy in the fellowship of your Father and of the angels. My cup runs over with the heavenly sweetness of your word.

"To whom should we go?"

The cords of death encompass me; a thousand dangers surround me; the enemies lie in wait for my soul; and I am helpless as a child and frail as a dry branch. But you are the Son of the living God, you are strong, you guard and preserve me, and in the midst of my weakness your power is made perfect; I am well content in weaknesses, in insults, in distresses, in anguishes for your sake, Lord Jesus; for when I am weak, then I am strong.

"To whom then should we go?"

You are the one who proclaims the Gospel to the poor, who had compassion on the sinful woman, whom men cast away, and who restored fallen Peter to grace. O Lord, how often I too have sinned, how frequently my heart has become lukewarm and cold toward you, who gave your life for me, and how my sins have gone over my soul like a flood. Ah, I had to hide myself like Adam from the righteous God and was near to despair. But you came to meet me as a merciful brother through this your word of life, you took me by the hand and led me again to Golgotha, where my sins were nailed to the cross, and led me, poor trembling child, to your Father, and he did not cast me out; for your sake he received me into grace again, and in blessed joy I again came to believe and know that you are the Son of the living God, and that there is salvation in none but you.

"To whom then should we go?"

In the bitter hour of death, in the dark passage through death’s valley, in the anguish of the last judgment—then you are my rod and staff; who is he that condemns? It is Christ who died, yes, who was raised and sits at the Father’s right hand; what shall I fear? "For this is the will of him who sent Jesus, that every one who sees the Son and believes in him should have eternal life, and I will raise him up on the last day," says the Lord. So that I am sure that nothing can separate me from the love of God in Jesus Christ, if only I am preserved in this confident confession:

"You have the words of eternal life, and I have believed and known that you are the Christ, the Son of the living God."

"To whom then should we go?"

\textbf{But ah—one of the twelve betrayed him.}

Is it I? Can it be I? the disciples asked in dread at the last supper.

Is it I? Is it I? We too must often ask in dread.

Therefore, "be not proud, but fear!"

